% SPDX-License-Identifier: MIT
\chapter{The interface}
\label{ch:interface}

\section{Layout of the window}

The window is a canvas surrounded by dockable panels. Every panel can be shown or
hidden from the \textbf{View} menu, and each has a shortcut
(Appendix~\ref{ch:shortcuts}).

\begin{center}
\begin{tabular}{@{}llp{0.52\linewidth}@{}}
\toprule
\textbf{Panel} & \textbf{Key} & \textbf{Purpose} \\
\midrule
Tree      & \keys{Alt+1} & The tree as a collapsible outline. Selecting a node here selects it on the canvas, and the reverse. \\
Style     & \keys{Alt+2} & Layout mode, branch mode, spacing, labels, theme. Most controls act immediately. \\
Tracks    & \keys{Alt+3} & The annotation stack. Reorder, hide, edit or remove tracks. \\
Search    & \keys{Ctrl+F} & Find taxa by name, with regular-expression support; matches are highlighted and selectable in bulk. \\
Inspector & \keys{Alt+4} & Everything known about the current selection: name, branch length, support, depth and any parsed annotations. \\
\bottomrule
\end{tabular}
\end{center}

\section{The canvas}

\subsection{Navigating}

\begin{itemize}
  \item \textbf{Zoom} with the mouse wheel. Zoom is anchored at the pointer, and
        it never re-runs the layout --- so zooming is instant regardless of tree
        size, and nothing shifts under the cursor.
  \item \textbf{Pan} by dragging with the middle button, or by dragging on empty
        canvas.
  \item \keys{Ctrl+0} fits the whole figure to the window;
        \keys{Ctrl+Shift+0} returns to actual size.
\end{itemize}

\subsection{Selecting}

Click a node to select it. Drag on empty canvas to rubber-band select a region.
\keys{Ctrl+Shift+A} extends the selection to the whole clade below the current
node, \keys{Ctrl+A} selects everything, and \keys{Esc} clears the selection.

Selection drives the tree operations in Chapter~\ref{ch:editing}: rerooting,
collapsing and pruning all act on what is selected.

\subsection{Context menu}

Right-clicking a node offers the operations that apply to it --- reroot here,
collapse, extract subtree, remove --- without going to the menu bar.

\section{The command palette}

\keys{Ctrl+P} opens a searchable list of every command in the application. It is
usually the fastest way to reach something whose menu location you cannot
remember, and it shows each command's keyboard shortcut alongside it, so it
doubles as a way to learn them.

The palette, the menus and the toolbar are all views of one registry of actions.
An action therefore cannot appear in one and be missing from another, and a
shortcut you rebind changes everywhere at once.

\section{Undo}

Every edit goes through a single command pipeline, so everything is undoable with
\keys{Ctrl+Z} and redoable with \keys{Ctrl+Shift+Z}. This includes operations
performed from the panels, not only those from the canvas.

Dragging a slider coalesces into a single history entry rather than one per pixel
of movement, so undo steps back over the whole drag.

\section{Themes}

\keys{Ctrl+T} switches between the light and dark themes. The theme belongs to
the figure, not merely to the screen: the composed scene carries its own
background colour, so an export made under the dark theme is dark in the file
too.

\galleryfig{feature-dark-theme}{The dark theme. Exports match the screen because
the background is part of the scene rather than a property of the window.}{0.9}

\section{Diagnostics}

When a file is loaded, anything \app{} had to interpret is recorded as a
diagnostic rather than discarded --- a duplicate tip name, a negative branch
length, an internal label that could be read either as a name or as a support
value. The status bar reports how many there were, and \keys{Ctrl+D} opens the
full list.

These are informational far more often than they are errors. A clean file
frequently produces one or two, and Chapter~\ref{ch:formats} explains the common
ones.
